\documentclass[aspectratio=169,10pt]{beamer}

\usepackage{booktabs}
\usepackage{graphicx}
\usepackage{amsmath}
\usepackage{xcolor}

\usetheme{Madrid}
\usecolortheme{seahorse}
\setbeamertemplate{navigation symbols}{}
\setbeamertemplate{caption}[numbered]
\setbeamerfont{title}{size=\LARGE,series=\bfseries}
\setbeamerfont{frametitle}{size=\large,series=\bfseries}
\setbeamertemplate{itemize items}[circle]

\title{When does A* save work?}
\subtitle{T01 checkpoint: question, contracts and frozen protocol}
\author[E. Haqqverdiyev, R. Iskandarli]{Elnur Haqqverdiyev \and Rajab Iskandarli}
\institute[UFAZ]{UFAZ, M1 Data Science and AI\\Advanced Data Structures and Algorithms, Fall 2026}
\date{Checkpoint}

\begin{document}

\begin{frame}
  \titlepage
\end{frame}

\begin{frame}{The question is not ``is A* faster?''}
  \begin{columns}[T,onlytextwidth]
    \begin{column}{0.55\textwidth}
      \textbf{Baseline:} Dijkstra with a binary heap and lazy deletion.\\
      \textbf{Comparison:} A* with the same heap and loop, plus a heuristic.
      \vspace{0.4em}

      A* can reduce \emph{expanded vertices}, but it pays for:
      \begin{itemize}
        \item one heuristic evaluation per generated vertex;
        \item larger heap keys and more memory traffic;
        \item extra work that is wasted if the heuristic is uninformative.
      \end{itemize}
    \end{column}
    \begin{column}{0.42\textwidth}
      \begin{block}{Research question}
        How do graph structure and heuristic strength change expanded
        vertices and runtime of A* versus Dijkstra, on unit-cost grids and a
        real road graph?
      \end{block}
    \end{column}
  \end{columns}
  \vspace{0.4em}
  \centering
  \textbf{Deliverable:} a rule for when the expansion saving pays for the
  guidance cost.
\end{frame}

\begin{frame}{Predictions to test}
  \small
  \setlength{\tabcolsep}{4pt}
  \begin{tabular}{p{0.05\textwidth}p{0.78\textwidth}p{0.11\textwidth}}
    \toprule
    & prediction & \\
    \midrule
    H1 & open grid: A* saves many expansions and time & \\
    H2 & random obstacles: saves, but less than open & \\
    H3 & clustered/rooms: weaker; rooms weaker than random & \\
    H4 & maze: no expansion saving, A* may be slower & \\
    H5 & long queries benefit more than short & \\
    H6 & real ALT: strong at K=8; K=1 may not help; diminishing at K=16 & \\
    H7 & \texttt{larger\_g} reduces median expansions on open grid & \\
    \bottomrule
  \end{tabular}
  \vspace{0.8em}

  These are written in \texttt{HYPOTHESES.md} before the main run; a negative
  result is acceptable.
\end{frame}

\begin{frame}{Data choice}
  \begin{columns}[T,onlytextwidth]
    \begin{column}{0.48\textwidth}
      \textbf{Grids as controls}
      \begin{itemize}
        \item open, random 20\%, clustered 30\%
        \item rooms with doors
        \item perfect maze
        \item 128, 256, 512, 1024 square
        \item 4-neighbour, unit cost
      \end{itemize}
      Why: isolate one structural property at a time.
    \end{column}
    \begin{column}{0.49\textwidth}
      \textbf{Real graph}
      \begin{itemize}
        \item SNAP \texttt{roadNet-CA}, California roads
        \item 1,957,027 vertices, 2,760,388 undirected edges after preparation
        \item largest connected component
        \item unit hop weight (SNAP has no lengths/travel times)
      \end{itemize}
      Why: sparse, planar-like, and exactly the setting where A* is usually
      recommended.
    \end{column}
  \end{columns}
\end{frame}

\begin{frame}{ADT contracts}
  \begin{columns}[T,onlytextwidth]
    \begin{column}{0.49\textwidth}
      \begin{block}{Graph: CSR}
        \texttt{offsets[n+1]}, \texttt{targets[m]}, \texttt{weights[m]}\\
        memory $O(n+m)$; degree $O(1)$; neighbours $O(\deg u)$
      \end{block}
      \begin{block}{Binary heap}
        push/pop $O(\log n)$; top $O(1)$\\
        duplicates allowed; stale entries discarded lazily\\
        tie-break: $(f,-g,id)$ or $(f,id)$
      \end{block}
    \end{column}
    \begin{column}{0.49\textwidth}
      \begin{block}{Search}
        one loop; Dijkstra = $h\equiv0$\\
        A* = $f=g+h(v)$; stop when target settled\\
        consistent $h$ $\Rightarrow$ no reopening\\
        O(1) scratch reset by generation stamp
      \end{block}
      \begin{block}{Heuristics}
        grids: Manhattan, admissible + consistent\\
        real: ALT $\max_L |d(L,v)-d(L,t)|$\\
        landmark BFS precomputed outside timing
      \end{block}
    \end{column}
  \end{columns}
\end{frame}

\begin{frame}{Correctness before benchmarking}
  \begin{columns}[T,onlytextwidth]
    \begin{column}{0.54\textwidth}
      \begin{itemize}
        \item hand-computed 5x5 wall grid (distance 8)
        \item six-vertex weighted graph vs Floyd--Warshall
        \item all five grid families vs independent BFS
        \item ALT admissibility/consistency for every free vertex
        \item path validation: contiguous edges and exact weight sum
        \item \texttt{expanded}, \texttt{stale\_popped}, \texttt{pushed} identities
      \end{itemize}
    \end{column}
    \begin{column}{0.43\textwidth}
      \begin{block}{Main-run guard}
        Every benchmark query is checked against BFS.  A distance mismatch
        aborts the run.  Timeouts are recorded, never silently dropped.
      \end{block}
      \begin{block}{Counter definitions}
        expanded = settled pop (target included)\\
        pushed = source + decreases\\
        stale = old entry or already closed
      \end{block}
    \end{column}
  \end{columns}
\end{frame}

\begin{frame}{Frozen protocol}
  \small
  \begin{columns}[T,onlytextwidth]
    \begin{column}{0.55\textwidth}
      \begin{tabular}{lcc}
        \toprule
        experiment & queries & timed runs\\
        \midrule
        grids (20 graphs) & 300 & 2,700\\
        real graph & 30 & 900\\
        \midrule
        total & 330 & 3,600\\
        \bottomrule
      \end{tabular}
      \vspace{0.7em}

      3 grid repeats; 5 real repeats.  One untimed warm-up per treatment.
    \end{column}
    \begin{column}{0.43\textwidth}
      \begin{itemize}
        \item paired same query, alternating order
        \item median per query; raw repeats kept
        \item bootstrap 95\% CI + Wilcoxon
        \item 30 s/search, 4 GiB/process
        \item one CPU via \texttt{taskset}
        \item seed 20261030
      \end{itemize}
    \end{column}
  \end{columns}
  \vspace{0.5em}
  \centering
  Protocol fixed at this checkpoint; changes require a deviation note.
\end{frame}

\begin{frame}{Prototype status and next steps}
  \begin{columns}[T,onlytextwidth]
    \begin{column}{0.50\textwidth}
      \textbf{Working now}
      \begin{itemize}
        \item C++20 CSR graph, binary heap, Dijkstra/A*
        \item five grid generators
        \item Manhattan and ALT heuristics
        \item BFS/Floyd--Warshall oracles
        \item correctness test suite
        \item raw-CSV benchmark runner
      \end{itemize}
    \end{column}
    \begin{column}{0.47\textwidth}
      \textbf{Next}
      \begin{itemize}
        \item pilot run to confirm the time budget
        \item main grid and real runs
        \item paired analysis, CIs and figures
        \item defence: answer the question, including negative results
      \end{itemize}
    \end{column}
  \end{columns}
  \vspace{0.7em}
  \centering
  \large The implementation is the same loop for both treatments; only the
  heuristic changes.
\end{frame}

\end{document}
